\section{Lecture 23: Berkovich Spaces as Analytic Stacks}
\label{lec:23}

Clausen returns to the Berkovich
world of \cref{lec:22} and promotes the Berkovich spectrum of a Banach ring to
an honest analytic stack $\operatorname{Spec}^{\mathrm{Berk}}(R,\lvert\cdot\rvert)$,
sitting over the classical topological space $\MBerk(R,\lvert\cdot\rvert)$ and
carrying a structure sheaf of $\infty$-categories -- a theory of
quasi-coherent sheaves on the Berkovich space. The stack is cut out inside
$\mathcal N\times\operatorname{Spec}(R,\mathrm{triv})$ (norms $\times$ trivially
structured $R$) by the condition that the tautological multiplicative norm be
bounded by the given Banach norm. The heart of the lecture is a theorem: under
a mild hypothesis $(\star)$ (existence of a topologically nilpotent unit whose
norm strictly multiplies), $\operatorname{Spec}^{\mathrm{Berk}}$ is
\emph{affine}, corresponds to a \emph{gaseous} analytic ring structure on the
condensed ring $R$ depending only on $R$ (not on the norm), and its tautological
norm is the expected one, given by overconvergent functions on discs. The proof
runs through a trace-class calculational trick and Serre duality on
$\mathbb P^1$. A short closing section globalises to a notion of Berkovich
analytic space.

Throughout, ``ring'' means commutative ring and everything is (light)
condensed.

\subsection{Recollection: Berkovich spaces}
\label{subsec:23:recollection}

We recall the classical setup (\cref{def:22:banach-ring,def:22:berkovich-spectrum}).
A \emph{Banach ring} $(R,\lvert\cdot\rvert)$ is a ring $R$ with a norm
$\lvert\cdot\rvert\colon R\to\mathbb R_{\ge0}$ that is submultiplicative and
satisfies the triangle inequality,
\begin{equation*}
\lvert xy\rvert\le\lvert x\rvert\,\lvert y\rvert,\qquad
\lvert x+y\rvert\le\lvert x\rvert+\lvert y\rvert,\qquad
\lvert-x\rvert=\lvert x\rvert,\qquad
\lvert1\rvert=1\ \text{ (unless } R=0\text{)},
\end{equation*}
with respect to which $R$ is complete.

To $(R,\lvert\cdot\rvert)$ Berkovich attaches a compact Hausdorff space, a subset
of a product of intervals,
\begin{equation*}
\MBerk(R,\lvert\cdot\rvert)\ \subseteq\ \prod_{f\in R}\,[0,\lvert f\rvert],
\end{equation*}
consisting of the \emph{multiplicative} seminorms bounded by
$\lvert\cdot\rvert$. A point is denoted $x$, but literally it is a valuation
$\lvert\cdot\rvert_x\colon R\to\mathbb R_{\ge0}$ which is strictly
multiplicative and satisfies the triangle inequality,
\begin{equation*}
\lvert fg\rvert_x=\lvert f\rvert_x\,\lvert g\rvert_x,\qquad
\lvert f+g\rvert_x\le\lvert f\rvert_x+\lvert g\rvert_x,\qquad
\lvert1\rvert_x=1 .
\end{equation*}
In general $R$ is not complete for a given $\lvert\cdot\rvert_x$; indeed there
can be many elements of norm $0$.

\begin{remark}[Residue fields and the three cases]
\label{rmk:23:residue-fields}
Completing $R$ with respect to a point $\lvert\cdot\rvert_x$ produces a
complete valued field
\begin{equation*}
\widehat{\kappa(x)},
\end{equation*}
the \emph{residue field} of $x$ in Berkovich's sense. There are essentially
three cases.
\begin{enumerate}
\item $\widehat{\kappa(x)}$ is archimedean; by Ostrowski this means
$(\mathbb R,\lvert\cdot\rvert^\alpha)$ or $(\mathbb C,\lvert\cdot\rvert^\alpha)$
with $\alpha\in(0,1]$. In the archimedean case there is no variety among
complete normed fields, only the choice of exponent.
\item $\lvert\cdot\rvert_x$ is non-archimedean but $\kappa(x)$ is discrete; the
extreme case is the trivial norm $\lvert\cdot\rvert_0$.
\item $\widehat{\kappa(x)}$ is a genuine non-archimedean field with non-discrete
induced topology, carrying a normalised absolute value
$\lvert\cdot\rvert^\alpha$, $\alpha\in(0,\infty)$ (no preferred normalisation
absent a choice of pseudo-uniformizer and a real number).
\end{enumerate}
These are the possible residue-field behaviours one sees.
\end{remark}

\subsection{The Berkovich spectrum as an analytic stack}
\label{subsec:23:stack}

The goal is to promote $\MBerk(R,\lvert\cdot\rvert)$ to an analytic stack, using
the \emph{stack of norms} $\mathcal N$ of \cref{lec:20,lec:22}. Recall
(\cref{def:20:norm,not:22:stack-of-norms}) that a map from an analytic stack
$X$ to $\mathcal N$ is the same as a map
\begin{equation*}
N\colon\ \mathbb P^1_X\ \longrightarrow\ [0,\infty]
\end{equation*}
satisfying the norm axioms -- the most important being multiplicativity, which
must be phrased with some care but in the end is just multiplicativity.

We also investigated the geometry of $\mathcal N$ last time: it lies over the
\emph{extended} Berkovich spectrum of the integers (\cref{thm:22:image},
\cref{fig:22:MZ}). On the classical Berkovich spectrum one imposes the usual
triangle inequality, which is exactly what prevents taking an arbitrary power
$\alpha$ of the archimedean absolute value; dropping it, an arbitrary power
behaves like a norm but really only a \emph{quasi}-norm (a constant appears in
front of the triangle inequality). In $\mathcal N$ no triangle inequality is
imposed, and one obtains the non-strict version in which all powers of the usual
archimedean norm are allowed, together with a strange limit point where the norm
takes the value $\infty$ on natural numbers. Along the interior of the
archimedean segment one is seeing the gaseous $\mathbb R$-theory; along the
interior of a $p$-adic segment the gaseous $p$-adic numbers, with the norm
changing but the underlying analytic ring not; the outer ends carry
$\mathbb F_p$ and other characteristic-$p$ phenomena, while the archimedean end
carries characteristic $0$. Morally the points of $\mathcal N$ correspond to
(minimal choices of) complete valued fields -- reals, $p$-adics, discrete
$\mathbb F_p$, discrete $\mathbb Q$ -- a substitute for the notion of
multiplicative valuation. But one still has to feed in the Banach ring $R$ to
get something non-trivial.

\begin{definition}[Berkovich spectrum]
\label{def:23:berkovich-spectrum}
For a Banach ring $(R,\lvert\cdot\rvert)$, let
\begin{equation*}
\operatorname{Spec}^{\mathrm{Berk}}(R,\lvert\cdot\rvert)\ \subseteq\
\mathcal N\times\operatorname{Spec}(R,\mathrm{triv})
\end{equation*}
be the analytic stack whose $X$-points are pairs
\begin{equation*}
\bigl(\,N\colon\mathbb P^1_X\to[0,\infty],\quad R\to\mathcal O(X)\,\bigr)
\end{equation*}
-- a norm $N$ on $X$ together with a map of condensed rings from $R$ to the
functions on $X$ -- such that for all $f\in R$ the composite
$N(f)\colon X\to[0,\infty]$ (evaluate the coordinate of $\mathbb P^1$ at the
section of $\mathbb A^1$ determined by $f$, then apply $N$) lands inside the
subspace $[0,\lvert f\rvert]$:
\begin{equation}
\label{eq:23:boundedness}
N(f)\ \in\ [0,\lvert f\rvert]\qquad\forall f\in R .
\end{equation}
Here $\operatorname{Spec}(R,\mathrm{triv})$ is the affine analytic stack of the
Banach ring $R$, viewed as a light condensed ring via its topology and equipped
with the \emph{trivial} (maximal) analytic ring structure
$D(R,\mathrm{triv})=D^{\mathrm{cond}}(R)$: a map $X\to\operatorname{Spec}(R,\mathrm{triv})$
is exactly a map of condensed rings $R\to\mathcal O(X)$.
\end{definition}

Condition \eqref{eq:23:boundedness} is what ties the norm factor to the ring
factor.

\begin{theorem}
\label{thm:23:localizes}
$\operatorname{Spec}^{\mathrm{Berk}}(R,\lvert\cdot\rvert)$ is an analytic stack,
and it localizes along the Berkovich spectrum
$\MBerk(R,\lvert\cdot\rvert)$ in the sense of \cref{lec:22}: there is a map of
locales
\begin{equation}
\label{eq:23:loc-map}
\Loc^{\mathrm{op}}\bigl(\operatorname{Spec}^{\mathrm{Berk}}(R,\lvert\cdot\rvert)\bigr)
\ \longrightarrow\ \MBerk(R,\lvert\cdot\rvert)
\end{equation}
from the locale of open substacks -- monomorphisms into
$\operatorname{Spec}^{\mathrm{Berk}}$ that are shriekable and cohomologically
smooth -- to the opens of the compact Hausdorff space
$\MBerk(R,\lvert\cdot\rvert)$ in its usual topology.
\end{theorem}

\begin{remark}[Structure sheaf and comparison to Poineau]
\label{rmk:23:structure-sheaf}
In particular one obtains a structure sheaf -- indeed a structure sheaf of
$\infty$-categories, i.e.\ a theory of quasi-coherent sheaves -- on the usual
Berkovich space. When $\MBerk(R,\lvert\cdot\rvert)$ is metrizable and
finite-dimensional one can even ask for a genuine map of analytic stacks to it
(\cref{thm:20:compact-hausdorff}); this is basically all cases, and requires
only checking a connectivity property of the idempotent algebras attached to
closed subsets, which is a condition verifiable in practice. Clausen noted it
would be interesting to compare with Poineau, who described structure sheaves
in certain cases over the integers (e.g.\ the Berkovich line over $\mathbb Z$)
more or less by hand \cite{poineau2008ladroiteberkovichsur,poineau2012espacesberkovichsurz}.
The present approach instead \emph{defines} an object which is a priori a
structure sheaf and satisfies $\infty$-categorical descent automatically, at the
cost of then having to compute its values. For rings like the integers this
computation is quite feasible; for a more arbitrary Banach ring it is less
obvious. As in the Huber theory, extracting the ``correct'' rings of functions
on closed or open discs likely requires working in a derived sense, with derived
rings complete in the appropriate structure.
\end{remark}

\begin{remark}[Relation to the adic spectrum]
\label{rmk:23:spa-comparison}
When $R$ is Tate one also has the localization over Huber's $\Spa$
(\cref{lec:22}), and a commutative diagram in which the Huber space maps to
$\operatorname{Spec}^{\mathrm{Berk}}$ and the solid analytic ring maps to the
trivially-structured $R$. The Tate case is examined in more detail below.
\end{remark}

\subsubsection*{Charts}

To see that $\operatorname{Spec}^{\mathrm{Berk}}$ is an analytic stack (only an
accessibility statement is needed, since it is defined as a functor of points)
we exhibit charts. Recall (\cref{rmk:22:gaseous-cover}) the chart for the stack
of norms: the gaseous base ring gives a cover
\begin{equation}
\label{eq:23:norms-chart}
\operatorname{Spec}\bigl(\mathbb Z[\widehat q^{\pm1}]_\gasx\bigr)
\ \longrightarrow\ \mathcal N ,
\end{equation}
the universal locus inside $\mathbb P^1$ where $\lvert T\rvert=\tfrac12$
(an element of norm $\tfrac12$ has been adjoined). Pulling back
\eqref{eq:23:norms-chart} along the structure map
$\operatorname{Spec}^{\mathrm{Berk}}\to\mathcal N$ produces a chart
$Y$ (a universal annulus):

% Provenance: board 11 (00:25:14) -- the pullback square defining the chart Y.
\begin{equation*}
\begin{tikzcd}[column sep=large, row sep=large]
Y \arrow[r] \arrow[d] &
\operatorname{Spec}\bigl(\mathbb Z[\widehat q^{\pm1}]_\gasx\bigr)
  \arrow[d] \\
\operatorname{Spec}^{\mathrm{Berk}}(R,\lvert\cdot\rvert) \arrow[r] & \mathcal N .
\end{tikzcd}
\end{equation*}

\begin{proposition}[The chart is affine]
\label{prop:23:chart-affine}
The chart $Y$ is affine: concretely, one takes
$\operatorname{Spec}(\mathbb Z[\widehat q^{\pm1}]_\gasx)\times
\operatorname{Spec}(R,\mathrm{triv})$ and passes to the closed subsets given by
the idempotent algebras obtained as $N(f)^{-1}$ of the closed subsets
$[0,\lvert f\rvert]\subseteq[0,\infty]$. On underlying analytic rings
\begin{equation}
\label{eq:23:Y-ring}
Y\ =\ \bigl(R\otimes_{\mathbb Z}\mathbb Z[\widehat q^{\pm1}]_\gasx\bigr)^{\gasx}.
\end{equation}
\end{proposition}

For a completely general $R$ the difficult part of computing $Y$ is already the
first step: forming the tensor product \eqref{eq:23:Y-ring} and identifying the
resulting analytic ring, in particular its underlying condensed ring. This
involves computing the gaseous localization of $R$ and then applying the gaseous
localization once more (recall from \cref{lec:14} that the gaseous localization
is a sequential colimit of Koszul complexes). Once that ring is known, passing to
the idempotent algebras cutting out the discs is not the hard part.

\begin{remark}[Descent happens at the zeroth stage]
\label{rmk:23:descent-zeroth}
The descent for the cover $Y\to\operatorname{Spec}^{\mathrm{Berk}}$ is
straightforward. As we saw when checking that this map is a cover -- it is
proper and descendable -- the unit object (structure sheaf) downstairs is
already a direct \emph{retract} of the pushforward of the structure sheaf on
$Y$. So the structure sheaf on $\operatorname{Spec}^{\mathrm{Berk}}$ is read off
as the retract: passing from $\operatorname{Spec}^{\mathrm{Berk}}$ up to $Y$
adjoins an extra variable $q$ with certain analytic properties, and one simply
takes the possible zeroth $q$-coefficients of the ring computed on $Y$ (a linear
retraction, exactly as constant coefficients sit inside Laurent series). This is
a retraction of modules, not of rings.
\end{remark}

\subsubsection*{The map to the Berkovich spectrum}

On $\operatorname{Spec}^{\mathrm{Berk}}(R,\lvert\cdot\rvert)$ one has the
universal norm $N$, and by fiat, for every $f\in R$, a map to the structure
sheaf. Together these give
\begin{equation}
\label{eq:23:to-product}
\operatorname{Spec}^{\mathrm{Berk}}(R,\lvert\cdot\rvert)\ \longrightarrow\
\prod_{f\in R}\,[0,\lvert f\rvert].
\end{equation}
Because we enforced $N(f)\le\lvert f\rvert$ for every element
(\eqref{eq:23:boundedness}), in particular $N(2)\le2$, which by the universal
triangle inequalities of \cref{prop:22:triangle} forces $N$ to satisfy the
triangle inequality. Hence \eqref{eq:23:to-product} lands inside the Berkovich
spectrum $\MBerk(R,\lvert\cdot\rvert)\subseteq\prod_{f\in R}[0,\lvert f\rvert]$,
giving the underlying map of \eqref{eq:23:loc-map}.

\subsection{Affineness under a topologically nilpotent unit}
\label{subsec:23:affine}

In general $\operatorname{Spec}^{\mathrm{Berk}}$ is \emph{not} affine.

\begin{example}[The integers with the archimedean norm]
\label{ex:23:integers}
For general $(R,\lvert\cdot\rvert)$ the stack
$\operatorname{Spec}^{\mathrm{Berk}}(R,\lvert\cdot\rvert)$ is not affine. For
instance
\begin{equation*}
\operatorname{Spec}^{\mathrm{Berk}}(\mathbb Z,\lvert\cdot\rvert_\infty)
\ =\ \mathcal N_{\lvert2\rvert\le2},
\end{equation*}
the locus $\{\lvert2\rvert\le2\}$ inside the stack of norms
($\lvert\cdot\rvert_\infty$ is the maximal norm on $\mathbb Z$, so this gives the
biggest possible Berkovich spectrum). It really is a stack, as one sees at the
points living at the boundaries.
\end{example}

We isolate the hypothesis that repairs this.

\begin{definition}[Assumption $(\star)$]
\label{def:23:star}
$(R,\lvert\cdot\rvert)$ satisfies $(\star)$ if (assuming $R\ne0$) there exists
$\pi\in R$ with
\begin{equation*}
\lvert\pi\rvert<1,\qquad \pi\ \text{a unit},\qquad
\lvert\pi f\rvert=\lvert\pi\rvert\,\lvert f\rvert\ \ \forall f\in R .
\end{equation*}
The last condition (the norm strictly multiplies against $\pi$) is equivalent,
for $R\ne0$, to $\lvert\pi^{-1}\rvert=\lvert\pi\rvert^{-1}$, as one sees from
submultiplicativity.
\end{definition}

\begin{example}[When $(\star)$ holds]
\label{ex:23:star-examples}
\leavevmode
\begin{itemize}
\item Any non-discrete valued field admits such a norm.
\item If $(R,\lvert\cdot\rvert)$ satisfies $(\star)$ and $\varphi\colon R\to R'$
is a contractive homomorphism ($\lvert\varphi(f)\rvert_{R'}\le\lvert f\rvert_R$
for all $f\in R$), then $R'$ satisfies $(\star)$ with respect to $\varphi(\pi)$:
the image is again a unit of norm $<1$, and $\lvert\varphi(\pi)^{-1}\rvert
=\lvert\varphi(\pi)\rvert^{-1}$ is more evidently preserved.
\item Any Tate--Huber ring has a norm defining its topology and satisfying
$(\star)$, with $\pi$ a pseudo-uniformizer; some mixed-characteristic examples
also arise.
\end{itemize}
Note $(\star)$ fails for $\mathbb Z$ with $\lvert\cdot\rvert_\infty$
(\cref{ex:23:integers}).
\end{example}

\begin{theorem}
\label{thm:23:affine}
If $(\star)$ holds, then
$\operatorname{Spec}^{\mathrm{Berk}}(R,\lvert\cdot\rvert)$ is affine, and
corresponds to an analytic ring structure on the condensed ring $R$. Moreover
this analytic ring structure depends only on the condensed ring $R$, not on the
choice of norm satisfying $(\star)$.
\end{theorem}

\subsubsection*{Producing the gaseous structure}

\begin{proof}[Proof of affineness]
Take $\pi$ as in $(\star)$. It gives a ring map $\mathbb Z[q]\to R$,
$q\mapsto\pi$. Since $\lvert\pi\rvert<1$, $\pi$ is topologically nilpotent (its
powers tend to $0$), so the map factors through the free condensed ring on a
topologically nilpotent element $\mathbb Z[\widehat q]$; since $\pi$ is a unit it
factors further:
\begin{equation*}
\begin{tikzcd}[column sep=large, row sep=large]
\mathbb Z[q] \arrow[r, "q\mapsto\pi"] \arrow[d] & R \\
\mathbb Z[\widehat q^{\pm1}] \arrow[ur, "q\mapsto\pi"'] &
\end{tikzcd}
\end{equation*}
We must check that $R$ is $q$-gaseous. Recall that the gaseous structure
(\cref{lec:14}) is defined by requiring the map
\begin{equation*}
1-q\cdot\mathrm{shift}\colon\ \Null(R)\ \xrightarrow{\ \sim\ }\ \Null(R)
\end{equation*}
to be an isomorphism, where mapping the universal null sequence $P$ into the
Banach space $R$ produces precisely the space of null sequences $\Null(R)$ in
$R$, and $1-q\cdot\mathrm{shift}$ is the endomorphism induced by
$1-T\cdot q\colon P\to P$. Its inverse is visibly the summation: if $(f_n)$ is a
null sequence, one can form
\begin{equation*}
\sum_n f_n\,\pi^n\ \in\ R
\end{equation*}
(the limit of a Cauchy sequence, using $\lvert\pi\rvert<1$ and the triangle
inequality). Thus $R$ is $q$-gaseous, and we take the induced analytic ring
structure from $\mathbb Z[\widehat q^{\pm1}]_\gasx$. Write $R^\gasx$ for the
resulting analytic ring (underlying condensed ring $R$).

Now recall that on $\operatorname{Spec}(\mathbb Z[\widehat q^{\pm1}]_\gasx)$ there
is a universal norm $N$ with
\begin{equation*}
N(q)\ \in\ \{\lvert\pi\rvert\}
\end{equation*}
(a singleton subspace of $[0,\infty]$, not merely a value; here $R\ne0$, so
$0<\lvert\pi\rvert<1$). By functoriality of norms under pullback we obtain a norm
$N$ on $R^\gasx$, i.e.\ on
$\operatorname{Spec}(R^\gasx)$, with induced analytic ring structure. The
content of the theorem is the following.

\begin{claim}
\label{claim:23:represents}
The normed analytic ring $\bigl(R^\gasx,N\bigr)$ represents
$\operatorname{Spec}^{\mathrm{Berk}}(R,\lvert\cdot\rvert)$.
\end{claim}

Over $\operatorname{Spec}(R^\gasx)$ we have the map to
$\operatorname{Spec}(R,\mathrm{triv})$ and the norm $N$ with
$N(\pi)\in\{\lvert\pi\rvert\}$, and it is straightforward from the construction
that $\operatorname{Spec}(R^\gasx)$ is universal with respect to that structure.
The difference from $\operatorname{Spec}^{\mathrm{Berk}}$ is that the latter
imposes a condition on the norm of \emph{every} element, whereas $R^\gasx$ was
built imposing a condition only on the norm of $\pi$. So the claim reduces to
the equivalence, for a norm $N$ on $R^\gasx$,
\begin{equation}
\label{eq:23:equivalence}
N(\pi)\in\{\lvert\pi\rvert\}
\quad\Longleftrightarrow\quad
N(f)\in[0,\lvert f\rvert]\ \ \forall f\in R .
\end{equation}
One direction is easy: given the right-hand side, apply it to $\pi$ and to
$\pi^{-1}$ to recover $N(\pi)=\lvert\pi\rvert$. The substance is that pinning
down $N(\pi)$ already constrains the norm of every element; this we prove by
computing the universal norm on $\operatorname{Spec}(R^\gasx)$.
\end{proof}

\subsubsection*{Computing the universal norm}

A norm in our sense on an analytic ring implicitly specifies the
\emph{overconvergent} functions on the disc of arbitrary radius $c$ centred at
the origin. The claim is that these are the usual ones from Berkovich theory.

\begin{proposition}[Overconvergent functions on a disc]
\label{prop:23:overconvergent}
For $0<c$, the value of the norm on $\operatorname{Spec}(R^\gasx)$ is given by
\begin{equation}
\label{eq:23:overconvergent}
\mathcal O_R\bigl(\{\lvert T\rvert\le c\}\bigr)\ =\
\varinjlim_{r>c}\ \mathcal O_R^{\lvert T\rvert\le r},
\qquad
\mathcal O_R^{\lvert T\rvert\le r}\ =\ \Bigl\{\textstyle\sum_n f_n T^n\ :\
\sum_n\lvert f_n\rvert\,r^n<\infty\Bigr\},
\end{equation}
the filtered colimit (in condensed rings) over radii $r>c$ of the universal
Banach ring in which the norm is $\le r$, i.e.\ power series over $R$ whose
coefficients decay fast enough that the substitution $T\mapsto r$ converges.
\end{proposition}

To make the calculation, rather than expanding
$\mathbb Z[\widehat q^{\pm1}]_\gasx\to\mathbb Z[\widehat q^{\pm1}]_\gasx
\otimes^{\mathbb L}_{\mathbb Z}R$ naively as a filtered colimit of copies of
$P$ -- where it is not even easy to see concentration in degree $0$ -- we use a
purely formal trick.

\begin{lemma}[Trace-class colimit trick]
\label{lem:23:trace-class}
Let $(\mathcal C,\otimes)$ be a closed symmetric monoidal $\infty$-category, and
let
\begin{equation*}
\cdots\to X_3\to X_2\to X_1
\end{equation*}
be a tower in $\mathcal C$ in which each transition map is \emph{trace class} --
a map $X\to Y$ is trace class if it comes from a map
$\mathbf 1\to X^\vee\otimes Y$, where $X^\vee=\underline{\Homop}(X,\mathbf 1)$ is
the internal dual (no dualizability assumed). Then for all $Y\in\mathcal C$,
\begin{equation*}
\varinjlim_n\ X_n^\vee\otimes Y\ \xrightarrow{\ \sim\ }\
\varinjlim_n\ \underline{\Homop}(X_n,Y),
\end{equation*}
an equality of internal (ind-)objects; if $\mathcal C$ has colimits and $\otimes$
commutes with colimits -- as in our examples -- the quotation marks disappear
and this is a genuine colimit.
\end{lemma}

\begin{proof}[Idea]
Given a trace-class map $X\to Y$, tensoring the datum $\mathbf 1\to X^\vee\otimes
Y$ with $X$ and using evaluation $X\otimes X^\vee\to\mathbf 1$ produces the map
$X\to Y$; dually it produces backward maps between the two ind/pro-systems
raising the index by one step, and these witness the isomorphism.
\end{proof}

Applying \cref{lem:23:trace-class} with $Y=\mathbf 1$, one recognises the
structure-sheaf object as a colimit over the duals $P^\vee$ with trace-class
transition maps; this reduces the computation to null sequences in $R$ and a
filtered colimit of such (any two versions of the unit disc agree once made
overconvergent, by sandwiching). To identify the dual of the overconvergent disc
with functions on the open complement one uses Serre duality.

\begin{remark}[Serre duality on $\mathbb P^1$]
\label{rmk:23:serre-duality}
Working over the gaseous base (the choice is immaterial), algebraic Serre duality
in the six-functor formalism says that $\mathbb P^1\to\operatorname{Spec}$ is
smooth and proper with dualizing object $\omega^1[1]$. Taking the (internal)
linear dual of the overconvergent disc then identifies it with the ring of
functions on the open complement. The finiteness needed for the double duality
is arranged by an overconvergence trick: the overconvergent disc is an
ind-object, whose dual is the pro-object obtained by \emph{increasing} the
radius; viewing that pro-object as an inverse limit of overconvergent discs with
the non-strict inequalities and dualizing termwise gives back the
overconvergent disc as a colimit of duals of $P$'s -- exactly the shape needed
for \cref{lem:23:trace-class}.
\end{remark}

It remains to supply the trace-class hypothesis; it holds for soft reasons.

\begin{lemma}[Trace-class claim]
\label{lem:23:trace-class-topology}
Let $X$ be a topological space and $Z\subseteq Z'$ closed subsets such that
there is an open $U$ with
\begin{equation*}
Z\ \subseteq\ U\ \subseteq\ Z'
\end{equation*}
(two closed subsets separated by an open subset). Then the restriction map
between the pushforwards of the constant sheaves is trace class in the derived
category of sheaves on $X$ with values in $D(\text{base})$.
\end{lemma}

Here the coefficient category is $D$ of the base ring, not the derived category
of $X$. Applying this on the interval $[0,\infty]$ -- where a closed disc of one
radius sits inside an open disc inside a closed disc of larger radius -- every
transition map in \eqref{eq:23:overconvergent} is trace class; pulling back along
the symmetric monoidal functor $D(\mathrm{Shv}([0,\infty]))\to D(\mathbb P^1)$
gives a trace-class map on $\mathbb A^1$, and a further trick shows its image
under the forgetful functor to the base is trace class as well (using that
trace-class maps form a two-sided ideal). Hence the restriction maps between the
overconvergent discs are trace class, and \cref{prop:23:overconvergent} follows,
giving the presentation \eqref{eq:23:overconvergent} that lets one calculate.

\begin{proposition}[The universal norm is the expected one]
\label{prop:23:conclusion-norm}
The norm $N$ on $\operatorname{Spec}(R^\gasx)$ is given by the usual
overconvergent functions on discs over $R$. In particular, to prove the
non-trivial direction of \eqref{eq:23:equivalence} -- that
$N(f)\in[0,c]$ for all $f$ (with $c=\lvert f\rvert$) -- it suffices to show
\begin{equation}
\label{eq:23:quotient-is-R}
\mathcal O_R\bigl(\{\lvert T\rvert\le c\}\bigr)\big/(T-f)\ =\ R .
\end{equation}
\end{proposition}

\begin{proof}
The map is evaluation $T\mapsto f$. As Scholze observed, one does not even need
to analyse the short exact sequence of Banach spaces: mere existence of the map
suffices. Indeed $\mathcal O_R(\{\lvert T\rvert\le c\})$ is an idempotent algebra
over $\mathbb A^1$ (over the polynomial ring on one generator), so base-changing
along $R[T]/(T-f)=R$ makes it an idempotent algebra over $R$; and $R$ itself is an
idempotent algebra over $R$. Two idempotent algebras over $R$ are equal as soon
as there is an algebra map between them, and the unit map (evaluation $T\mapsto
f$) provides one. Hence \eqref{eq:23:quotient-is-R}.
\end{proof}

\begin{remark}[Why overconvergence is essential]
\label{rmk:23:overconvergence-essential}
The overconvergence is doing real work. In the non-archimedean (solid) setting
the affinoid algebra of the closed disc -- without overconvergence -- is
already idempotent, as we saw in the solid theory. In the archimedean case it is
not: the closed-disc algebra of convergent (rather than overconvergent) functions
is not idempotent (Weierstrass division only gives convergence on a smaller
disc), so the argument above would fail without passing to the overconvergent
colimit.
\end{remark}

\subsection{Independence of the norm}
\label{subsec:23:independence}

Two things remain in \cref{thm:23:affine}: that the analytic ring
$\operatorname{Spec}(R^\gasx)$ is independent of the choice of $\lvert\cdot\rvert$
and $\pi$, and (for the record) that the tautological norm on it is independent
of $\lvert\cdot\rvert$ up to rescaling. The first is achieved by an intrinsic
description of the gaseous structure.

\begin{proposition}[Intrinsic description]
\label{prop:23:intrinsic}
$R^\gasx$ is the initial analytic ring $A$ with a map
$\operatorname{Spec}(A)\to\operatorname{Spec}(R,\mathrm{triv})$ (i.e.\ a map
$R^{\mathrm{triv}}\to A$) such that for all topologically nilpotent units
$\pi\in R$, $A$ is $\pi$-gaseous -- i.e.\ the ring map
$\mathbb Z[q^{\pm1}]\to A^\triangleright$, $q\mapsto\pi$, yields a map of
analytic rings
\begin{equation*}
\mathbb Z[\widehat q^{\pm1}]_\gasx\ \longrightarrow\ A
\qquad(\text{equivalently } D(A)\subseteq\{\pi\text{-gaseous modules}\}).
\end{equation*}
This condition depends only on the topological ring $R$, not on the norm.
\end{proposition}

\begin{proof}
If $\pi$ is topologically nilpotent then there is $n>0$ with
$\lvert\pi^n\rvert<1$, and being $\pi$-gaseous is invariant under replacing $\pi$
by a power (a remark of Scholze's, made earlier); so we may assume
$0<\lvert\pi\rvert<1$. Then for the universal norm built over $R^\gasx$,
\begin{equation}
\label{eq:23:pi-in-01}
N(\pi)\ \in\ (0,1),
\end{equation}
which by the axioms of a normed analytic ring already implies $\pi$ is gaseous.
Now $R^\gasx$ was built to satisfy $(\star)$ only for some fixed $\pi_0$; the
computation \eqref{eq:23:pi-in-01} shows it satisfies the gaseous condition for
\emph{all} choices of $\pi$, giving the initial/universal property. Since the
condition never referred to the norm, $R^\gasx$ depends only on the condensed
ring $R$.
\end{proof}

\begin{proposition}[The norm, up to rescaling]
\label{prop:23:norm-rescaling}
The norm $N$ on $\operatorname{Spec}(R^\gasx)$ is independent of the choice of
$\lvert\cdot\rvert$ up to exponentiation by a map
\begin{equation*}
\alpha\colon\ \operatorname{Spec}(R^\gasx)\ \longrightarrow\ \mathbb R_{>0}.
\end{equation*}
\end{proposition}

\begin{proof}
Let $N$ come from $\lvert\cdot\rvert$ (satisfying $(\star)$ for $\pi$) and $N'$
from $\lvert\cdot\rvert'$ (satisfying $(\star)$ for $\pi'$, with
$0<\lvert\pi'\rvert'<1$). Then
\begin{equation}
\label{eq:23:Npiprime}
N(\pi')\colon\ \operatorname{Spec}(R^\gasx)\ \longrightarrow\ (0,1),
\end{equation}
and since $\mathbb R_{>0}$ acts simply transitively on $(0,1)$ by exponentiation
there is $\alpha$ with
\begin{equation}
\label{eq:23:Npiprime-alpha}
N(\pi')^{\alpha}\ =\ \{\lvert\pi'\rvert'\}.
\end{equation}
Then $N^{\alpha}$ and $N'$ are two norms on $\operatorname{Spec}(R^\gasx)$ with
the same (fixed) value on $\pi'$, so by the classification of norms
(\cref{thm:20:classification}) they are equal: $N'=N^\alpha$.
\end{proof}

\begin{remark}[Ofer's question]
\label{rmk:23:ofer}
Is the rescaling $\alpha\colon\operatorname{Spec}(R^\gasx)\to\mathbb R_{>0}$
pulled back from the Berkovich spectrum $\MBerk(R,\lvert\cdot\rvert)$ -- i.e.\
does it factor
\begin{equation*}
\begin{tikzcd}[column sep=large, row sep=large]
& \mathbb R_{>0} \\
\operatorname{Spec}(R^\gasx) \arrow[ur, "\alpha"] \arrow[r] &
\MBerk(R,\lvert\cdot\rvert) \arrow[u, dashed, "?"']
\end{tikzcd}
\end{equation*}
The answer is yes, and by construction: $\alpha$ is determined by $N(\pi')$
(\eqref{eq:23:Npiprime-alpha}), and the Berkovich spectrum records the norms of
\emph{all} elements of $R$, in particular that of $\pi'$. So $N(\pi')$, and hence
$\alpha$, is pulled back from $\MBerk(R,\lvert\cdot\rvert)$.
\end{remark}

\subsection{Globalization}
\label{subsec:23:globalization}

Globalizing is now trivial. We may simply make a definition.

\begin{definition}[Berkovich analytic space]
\label{def:23:berkovich-analytic-space}
A \emph{Berkovich analytic space} is a triple $(X,S,\pi)$ consisting of an
analytic stack $X$, a locally compact Hausdorff space $S$, and a map of locales
\begin{equation*}
\pi\colon\ \Loc^{\mathrm{op}}(X)\ \longrightarrow\ S ,
\end{equation*}
such that, locally on $S$ for the local section topology, it is isomorphic to
\begin{equation*}
\bigl(\operatorname{Spec}^{\mathrm{Berk}}(R,\lvert\cdot\rvert),\
\MBerk(R,\lvert\cdot\rvert),\ \pi\bigr)
\end{equation*}
for some Banach ring $(R,\lvert\cdot\rvert)$.
\end{definition}

The point is that working locally on $S$ automatically means working locally on
$X$: by definition of a map of locales, an open cover of $S$ pulls back to a
cover of $X$ -- indeed a cover in the Grothendieck topology defining analytic
stacks, even an open cover in the sense of the six-functor formalism.

\begin{notation}[Local section topology]
\label{not:23:local-section}
On the category of locally compact Hausdorff spaces, a family $\{X_i\to S\}$ is a
covering for the \emph{local section topology} if for every point of $S$ there is
an open neighbourhood, an index $i$, and a section of the pullback of $X_i$ to
that neighbourhood; the maps $X_i\to S$ themselves are arbitrary (not required to
be open inclusions). This is the same Grothendieck topology as the one generated
by the usual open covers -- every covering sieve equals the sieve generated by
some open cover of $S$ -- but it is convenient for expressing the sense in which
a locally compact Hausdorff space ``is locally'' compact Hausdorff: the affine
models $\MBerk(R,\lvert\cdot\rvert)$ are compact Hausdorff, and open subsets of,
say, $\mathbb R$ are usually not compact, so the naive formulation fails, whereas
covering by maps with local sections (compact neighbourhoods of each point) works.
\end{notation}

\begin{remark}[A derived fly in the ointment]
\label{rmk:23:derived-subtlety}
As with the Huber theory (\cref{lec:10}), the presentation starts from classical
objects -- Banach rings -- and when one works locally on the spectrum of such
an object one may a priori encounter derived objects. It is slightly more acute
here: the basic building blocks are Banach rings, but the structure sheaf is
overconvergent, so it is essentially never a sheaf of Banach rings, and there is
a mismatch between the global object one starts with and the local values. In
practice this does not obstruct anything -- the neighbourhood bases one computes
land in degree $0$ and fit back into the framework -- and the definitions can be
modified to accommodate the general case. For instance, one may enlarge the class
of affine models by passing to arbitrary inverse limits of the
$\operatorname{Spec}^{\mathrm{Berk}}(R,\lvert\cdot\rvert)$ (filtered colimits of
analytic rings, over inverse limits of compact Hausdorff spaces); the
overconvergent objects then count as affine, replacing Banach rings by condensed
rings with suitable norm data. We did not try to give the best possible
formulation, only to connect with the classical picture.
\end{remark}
